% ---
% title: Waiting for the bus
% tags: continuous, density, uniform, expectation, variance
% ---

\begin{problem}
	A bus arrives at a stop at a random time between $0$ and $15$ minutes from
	now. The arrival time is uniformly distributed on this interval: its
	density is constant on $[0,15]$ and $0$ outside. Let $X$ be your
	waiting time (in minutes).
	(a) Find the density $f(x)$ of $X$.
	(b) Find $P(X<5)$ and $P(3<X<8)$.
	(c) Find the expected waiting time $\E[X]$ and the standard deviation
	$\SD(X)$.
	(d) You have already waited $5$ minutes and the bus has not come yet. What
	is the probability that you wait more than $10$ minutes in total?
\end{problem}

\begin{slideproblem}
	A bus arrives at a random time between $0$ and $15$ minutes from now,
	uniformly distributed on $[0,15]$. Let $X$ be your waiting time
	(minutes).
	(a) Find the density $f(x)$.
	\ (b) $P(X<5)$ and $P(3<X<8)$.
	\ (c) $\E[X]$ and $\SD(X)$.
	\ (d) You have already waited $5$ minutes. Probability that you wait more
	than $10$ minutes in total?
\end{slideproblem}

\begin{slides}
	\small
	\exstatement
	\pause
	\begin{columns}[T]
		\begin{column}{.6\textwidth}
			\sol (a) Rectangle of width 15 and area 1:
			\ $\alert{f(x)=\frac1{15}}$ on $[0,15]$, $0$ outside.\\[1mm]
			\pause
			(b) Area $=$ width $\times$ height:
			$P(X<5)=\frac5{15}=\alert{\frac13}$, \
			$P(3<X<8)=\frac{5}{15}=\alert{\frac13}$.\\[1mm]
			\pause
			(c) Uniform on $(0,15)$: \ $\E[X]=\frac{0+15}{2}=\alert{7.5}$,\\
			$\Var(X)=\frac{15^2}{12}=18.75$, \ $\SD(X)\approx\alert{4.33}$.\\[1mm]
			\pause
			(d) $P(X>10\mid X>5)=\frac{P(X>10)}{P(X>5)}
				=\frac{5/15}{10/15}=\alert{\frac12}$,\\
			not $P(X>10)=\frac13$.
		\end{column}
		\begin{column}{.38\textwidth}
			\centering
			\begin{tikzpicture}[x=.26cm,y=28cm]
				\fill[teal!22] (3,0) rectangle (8,1/15);
				\draw[teal] (3,0)--(3,1/15) (8,0)--(8,1/15);
				\draw[->] (0,0)--(17,0) node[right,font=\scriptsize] {$x$};
				\draw[->] (0,0)--(0,.095) node[above,font=\scriptsize] {$f(x)$};
				\draw[navy,thick] (0,1/15)--(15,1/15);
				\draw[navy,dashed] (15,0)--(15,1/15);
				\draw[gold,thick,dashed] (7.5,0)--(7.5,.085);
				\node[above,font=\tiny,text=gold!80!black] at (7.5,.085) {$\E[X]=7.5$};
				\node[left,font=\tiny] at (0,1/15) {$\frac1{15}$};
				\node[font=\tiny,text=teal!70!black] at (5.5,.033) {$\frac13$};
				\foreach \x in {3,8,15}
				\node[below,font=\tiny] at (\x,0) {$\x$};
			\end{tikzpicture}
		\end{column}
	\end{columns}
\end{slides}

\begin{solution}
	\idea The density is a constant $c$ on $[0,15]$, and the total area under a
	density is $1$; this gives $c$. A probability is the area under the density
	over the interval. Here that area is a rectangle, so no integral is really
	needed: the probability is width times height. A random variable with a
	constant density on an interval $(\alpha,\beta)$ is called
	\emph{uniform} on $(\alpha,\beta)$; its expected value is the midpoint
	$\frac{\alpha+\beta}{2}$ and its variance is $\frac{(\beta-\alpha)^2}{12}$.

	\step{(a)} Write $f(x)=c$ for $0\le x\le 15$ and $f(x)=0$ otherwise. The
	total probability is $1$, so
	\[
		1=\int_{-\infty}^{\infty}f(x)\,dx=\int_{0}^{15}c\,dx
		=\Big[c\,x\Big]_{0}^{15}=15c-0=15c .
	\]
	Hence $c=\frac1{15}$:
	\[
		f(x)=\begin{cases}
			\dfrac1{15} & 0\le x\le 15, \\[6pt]
			0           & \text{otherwise.}
		\end{cases}
	\]
	Geometrically: the graph of $f$ is a rectangle of width $15$, and its area
	must be $1$, so its height is $\frac1{15}$.
	\begin{center}
		\begin{tikzpicture}[x=0.55cm,y=30cm,>=stealth]
			\fill[teal!25] (3,0) rectangle (8,1/15);
			\fill[gold!40] (0,0) rectangle (3,1/15);
			\draw[->] (-1,0)--(17,0) node[right] {$x$};
			\draw[->] (0,0)--(0,0.1) node[above] {$f(x)$};
			\draw[navy,very thick] (-1,0)--(0,0);
			\draw[navy,very thick] (0,1/15)--(15,1/15);
			\draw[navy,very thick] (15,0)--(16.8,0);
			\draw[navy,densely dotted] (15,0)--(15,1/15);
			\node[left] at (0,1/15) {\footnotesize$\frac1{15}$};
			\foreach \x in {3,5,8,10,15}
			\draw (\x,0.002)--(\x,-0.002) node[below] {\footnotesize$\x$};
			\draw[navy,thick,dashed] (5,0)--(5,1/15);
			\node[below left] at (0,0) {\footnotesize$0$};
		\end{tikzpicture}
	\end{center}
	The gold strip plus the left part of the teal strip (up to the dashed line
	at $5$) is the area $P(X<5)$; the teal strip from $3$ to $8$ is the area
	$P(3<X<8)$.

	\step{(b)} $X<5$ means the bus comes within the first $5$ minutes. The
	probability is the area under $f$ between $0$ and $5$ (there is no
	probability below $0$):
	\[
		P(X<5)=\int_{0}^{5}\frac1{15}\,dx=\Big[\frac{x}{15}\Big]_{0}^{5}
		=\frac5{15}-0=\frac13 .
	\]
	The same with geometry: a rectangle of width $5$ and height $\frac1{15}$.

	Next, again the area under $f$, now between $3$ and $8$:
	\[
		P(3<X<8)=\int_{3}^{8}\frac1{15}\,dx=\Big[\frac{x}{15}\Big]_{3}^{8}
		=\frac8{15}-\frac3{15}=\frac5{15}=\frac13 .
	\]
	This equals $P(X<5)$. The reason: for a uniform random variable
	the density is flat, so the probability of an interval inside $[0,15]$ is
	its length divided by $15$. Only the length of the interval matters, not
	where it lies. Both $(0,5)$ and $(3,8)$ have length $5$.

	\step{(c)} The expected value is $\E[X]=\int x\,f(x)\,dx$, and the
	density is $\frac1{15}$ on $[0,15]$:
	\[
		\E[X]=\int_0^{15}\frac{x}{15}\,dx=\Big[\frac{x^2}{30}\Big]_0^{15}
		=\frac{225}{30}=7.5 .
	\]
	This is the midpoint of the interval, as the symmetry of the rectangle
	suggests. For the variance, first
	\[
		\E[X^2]=\int_0^{15}\frac{x^2}{15}\,dx=\Big[\frac{x^3}{45}\Big]_0^{15}
		=\frac{3375}{45}=75,
	\]
	then
	\[
		\Var(X)=\E[X^2]-(\E[X])^2=75-7.5^2=75-56.25=18.75,\qquad
		\SD(X)=\sqrt{18.75}\approx4.33 .
	\]
	The general formulas for the uniform distribution on $(\alpha,\beta)$
	give the same: $\E[X]=\frac{0+15}{2}=7.5$ and
	$\Var(X)=\frac{(15-0)^2}{12}=\frac{225}{12}=18.75$. On average you wait
	$7.5$ minutes, and the waiting time typically differs from this by about
	$4.3$ minutes.

	\step{(d)} This is a conditional probability. Let $E=\{X>10\}$ (you wait
	more than $10$ minutes in total) and $F=\{X>5\}$ (the bus has not come in
	the first $5$ minutes). By definition
	\[
		P(E\mid F)=\frac{P(EF)}{P(F)} .
	\]
	If $X>10$, then surely $X>5$, so $EF=E$ and $P(EF)=P(X>10)$. Both
	probabilities are areas of rectangles of height $\frac1{15}$:
	\[
		P(X>10)=\int_{10}^{15}\frac1{15}\,dx=\frac{15-10}{15}=\frac5{15},
		\qquad
		P(X>5)=\int_{5}^{15}\frac1{15}\,dx=\frac{15-5}{15}=\frac{10}{15}.
	\]
	Therefore
	\[
		P(X>10\mid X>5)=\frac{5/15}{10/15}=\frac5{10}=\frac12 .
	\]
	Interpretation: once you know $X>5$, the bus comes at a time uniformly
	spread over the remaining $10$ minutes $(5,15]$, and ``more than $10$ in
	total'' is the last $5$ of these $10$ minutes.

	\begin{answer}
		(a) ${f(x)=\frac1{15}}$ for ${0\le x\le15}$, and $0$ otherwise \quad
		(b) ${P(X<5)=\frac13}$, ${P(3<X<8)=\frac13}$ \quad
		(c) $\E[X]=7.5$ minutes, $\Var(X)=18.75$, $\SD(X)\approx4.33$ minutes \quad
		(d) ${P(X>10\mid X>5)=\frac12}$
	\end{answer}
	\begin{remark}
		\textbf{Common mistake:} in (d), answering $P(X>10)=\frac13$. The
		information ``you have already waited $5$ minutes'' changes the
		probability: the first $5$ minutes are ruled out, and the remaining
		$10$ minutes become the whole sample space. The conditional probability
		$\frac12$ is larger than the unconditional $\frac13$, because the $10$
		minutes that are left now carry all the probability.
		Note also that for a continuous random variable $P(X=a)=0$ for every
		$a$, so it does not matter whether we write $X<5$ or $X\le5$.
	\end{remark}
\end{solution}
